% -*- coding: utf-8 ; -*-
\documentclass[dvipsnames]{article}
\usepackage[french]{babel}
\frenchsetup{og = « , fg = »}

\usepackage[escape-inside=$$]{piton}

\usepackage{xcolor}

\usepackage{geometry}
\geometry{left=2.8cm,right=2.8cm,top=2.5cm,bottom=2.5cm,papersize={21cm,29.7cm}}

\usepackage{enumitem}
\usepackage{verbatim}
\usepackage{amsmath}
\usepackage{tabularx}
\usepackage{booktabs}

\usepackage{caption}
\captionsetup{labelfont = bf}

% We use \MakeShortVerb of shortvrb and not \DefineShortVerb of fancyvrb
% because we don't want the contents of short verbatim colored in gray
\usepackage{shortvrb}
\MakeShortVerb{\|}


\skip\footins = 2\bigskipamount

\usepackage{fancyvrb}
\fvset{commandchars=\~\#\@,formatcom=\color{gray}}
\def\emphase{\bgroup\color{RoyalPurple}\let\next=}


\usepackage{titlesec}
\titlespacing*{\section}{0pt}{6.5ex plus 1ex minus .2ex}{4.3ex plus .2ex}
\titlespacing*{\subsection}{0pt}{4.5ex plus 1ex minus .2ex}{2ex plus .2ex}

\def\interitem{\vspace{7mm plus 2 mm minus 3mm}}          


\usepackage{footnote}

\usepackage[hyperfootnotes = false]{hyperref}

\hypersetup
  {
    pdfinfo = 
      {
        Title = L’extension piton ,
        Subject = Une extension LaTeX ,
        Author = F. Pantigny 
      }
  } 


\NewDocumentEnvironment {scope} {} {} {}

\NewDocumentCommand {\pkg} {m} {\textsf{#1}}
\NewDocumentCommand {\cls} {m} {\textsf{#1}}

\setlength{\parindent}{0pt}


\begin{document}

\VerbatimFootnotes


\title{L'extension LaTeX \pkg{piton}\thanks{Ce document correspond à la 
version~\myfileversion\space de \pkg{piton}, à la date du~\myfiledate.}} 
\author{F. Pantigny \\ \texttt{fpantigny@wanadoo.fr}}

\maketitle

\begin{abstract}
L'extension \pkg{piton} propose des outils pour composer du code Python avec une coloration syntaxique
en utilisant la bibliothèque Lua LPEG. L'extension \pkg{piton} nécessite l'emploi de LuaLaTeX.
\end{abstract}

\section{Présentation}

L'extension \pkg{piton} utilise la librarie Lua nommée LPEG\footnote{LPEG est une librairie de capture de motifs
  (\emph{pattern-matching} en anglais) pour Lua, écrite en C, fondée sur les PEG (\emph{parsing expression
    grammars}): \url{http://www.inf.puc-rio.br/~roberto/lpeg/}} pour «parser» le code Python et le composer avec un
coloriage syntaxique. Comme elle utilise du code Lua, elle fonctionne uniquement avec |lualatex| (et ne va pas
fonctionner avec les autres moteurs de compilation LaTeX, que ce soit |latex|, |pdflatex| ou |xelatex|). Elle
n'utilise aucun programme extérieur et la compilation ne requiert donc pas |--shell-escape|. La compilation est
très rapide puisque tout le travail du parseur est fait par la librairie LPEG, écrite en C.

\bigskip
Voici un exemple de code Python composé avec l'environnement |{Piton}| proposé par \pkg{piton}.

\bigskip
\begin{Piton}
from math import pi

def arctan(x,n=10):
    """Compute the value of arctan(x)
       n is the number of terms if the sum"""
    if x < 0:
         return -arctan(-x) # recursive call
    elif x > 1: 
         return pi/2 - arctan(1/x) 
         ## (on a utilisé le fait que $\arctan(x)+\arctan(1/x)=\frac{\pi}{2}$ pour $x>0$)\footnote{Cet échappement vers LaTeX a été obtenu en débutant par \ttfamily\#\#.}
    else 
         s = 0
         for k in range(n):
              s += (-1)**k/(2*k+1)*x**(2*k+1)
         return s 
\end{Piton}



\section{Installation}

L'extension \pkg{piton} est composée de deux fichiers : |piton.sty| et |piton.lua| (le fichier LaTeX |piton.sty|
chargé par |\usepackage| va à son tour charger le fichier |piton.lua|). Les deux fichiers doivent être présents
dans un répertoire où LaTeX pourra les trouver. Le mieux est de les installer dans une arborescence |texmf|.


\section{Utilisation de l'extension}

Pour utiliser l'extension \pkg{piton}, l'utilisateur final a seulement à charger l'extension dans son document avec
la commande classique |\usepackage| et se souvenir que la compilation doit être faite avec LuaLaTeX. 

\medskip
L'extension \pkg{piton} fournit trois outils pour composer du code Python: la commande |\piton|, l'environnement
|{Piton}| et la commande |\PitonInputFile|.

\begin{itemize}
\item La commande |\piton| doit être utilisée pour composer de petits éléments de code à l'intérieur d'un
paragraphe. \emph{Attention} : Cette fonction prend son argument en mode \emph{verbatim} (comme la commande
\verb|\verb|) et, de ce fait, cette fonction ne peut pas être utilisée à l'intérieur d'un argument d'une autre
fonction (on peut néanmoins l'utiliser à l'intérieur d'un environnement).

\item L'environnement |{Piton}| doit être utilisé pour composer des codes de plusieurs lignes. 

\item La commande |\PitonInputFile| doit être utilisée pour insérer et composer un fichier extérieur.
\end{itemize}

\bigskip
Il est possible de composer des commentaires en LaTeX en commençant par |##| (c'est un échappement vers LaTeX). Les
caractères |##| eux-mêmes ne seront pas imprimés et les espaces qui les suivent sont supprimés.


\section{Personnalisation}

\subsection{La commande \textbackslash PitonOptions}

La commande |\PitonOptions| propose cinq clés : |gobble|, |auto-gobble|, |env-gobble|, |line-numbers| et |all-line-numbers|.

\begin{itemize}
\item La clé |gobble| peut comme valeur un entier positif $n$ : les $n$ premiers caractères de chaque ligne sont
alors retirés (avant formatage du code) dans les environnements |{Piton}|.

\item Quand la clé |auto-gobble| est activée, l'extension \pkg{piton} détermine la valeur minimale $n$ du nombre
d'espaces successifs débutant chaque ligne (non vide) de l'environnment |{Piton}| et applique |gobble| avec cette
valeur de~$n$.

* Quand la clé |env-gobble| est activée, \pkg{piton} applique la clé |gobble| avec une valeur de $n$ égale au
nombre d'espaces sur la ligne précédant le |\end{Piton}| (si cette ligne ne comporte que des espaces).

\item Avec la clé |line-numbers|, les lignes \emph{non vides} sont numérotées dans les environnements \verb|{Piton}|
et dans les listings produits par la commande |\PitonInputFile|.

\item Avec la clé |all-line-numbers|, \emph{toutes} les lignes sont numérotées, y compris les lignes vides.
\end{itemize}

\bigskip

\begingroup
\fvset{commandchars=\~\&\@,formatcom=\small\color{gray}}
\begin{Verbatim}
~emphase&\PitonOptions{line-numbers,auto-gobble}@
\begin{Piton}
    from math import pi

    def arctan(x,n=10):
        """Compute the value of arctan(x)
           n is the number of terms if the sum"""
        if x < 0:
             return -arctan(-x) # recursive call
        elif x > 1: 
             return pi/2 - arctan(1/x) 
             ## (on a utilisé le fait que $\arctan(x)+\arctan(1/x)=\frac{\pi}{2}$ pour $x>0$)
        else 
             s = 0
             for k in range(n):
                  s += (-1)**k/(2*k+1)*x**(2*k+1)
             return s 
\end{Piton}
\end{Verbatim}
\endgroup

\begingroup
\PitonOptions{line-numbers,auto-gobble}
\begin{Piton}
    from math import pi

    def arctan(x,n=10):
        """Compute the value of arctan(x)
           n is the number of terms if the sum"""
        if x < 0:
             return -arctan(-x) # recursive call
        elif x > 1: 
             return pi/2 - arctan(1/x) 
             ## (on a utilisé le fait que $\arctan(x)+\arctan(1/x)=\frac{\pi}{2}$ pour $x>0$)
        else 
             s = 0
             for k in range(n):
                  s += (-1)**k/(2*k+1)*x**(2*k+1)
             return s 
\end{Piton}
\endgroup

\subsection{L'option escape-inside}

L'option |escape-inside| doit être utilisée au chargement de \pkg{piton} (c'est-à-dire dans l'instruction
|\usepackage|). Pour des raisons techniques, elle ne peut pas être fixée dans |\PitonOptions|. Elle prend comme
valeur deux caractères qui seront utilisés pour délimiter des parties qui seront composées en LaTeX.


\bigskip
Dans l'exemple suivant, on suppose que l'extension \pkg{piton} a été chargée de la manière suivante :

\smallskip
\begin{Verbatim}
\usepackage[~emphase#escape-inside=$$@]{piton}
\end{Verbatim}

\medskip
Dans le code suivant, qui est une programmation récursive de la factorielle, on décide de surligner en jaune
l'instruction qui contient l'appel récursif.
\begin{Verbatim}
\begin{Piton}
def fact(n):
    if n==0:
        return 1
    else:
        ~emphase#$\colorbox{yellow!50}{$@return n*fact(n-1)~emphase#$}$@
\end{Piton}
\end{Verbatim}

\begin{Piton}
def fact(n):
    if n==0:
        return 1
    else:
        $\colorbox{yellow!50}{$return n*fact(n-1)$}$
\end{Piton}

\bigskip
\emph{Attention} : L'échappement vers LaTeX permis par les caractères de |escape-inside| n'est pas actif dans les
chaînes de caractères ni dans les commentaires (pour avoir un commentaire entièrement en échappement vers LaTeX,
il suffit de le faire débuter par |##|).

\subsection{Les styles}

L'extension \pkg{piton} fournit la commande |\SetPitonStyle| pour personnaliser les différents styles utilisés pour
formater les éléments syntaxiques des listings Python. Ces personnalisations ont une portée qui correspond au
groupe TeX courant.

\bigskip
La commande |\SetPitonStyle| prend en argument une liste de couples \textsl{clé=valeur}. Les clés sont les noms des
styles et les valeurs sont les instructions LaTeX de formatage correspondantes.

\bigskip
Ces instructions LaTeX doivent être des instructions de formatage du type de |\color{...}|, |\bseries|,
|\slshape|, etc. (les commandes de ce type sont parfois qualifiées de \emph{semi-globales}). Il est aussi possible
de mettre, \emph{à la fin de la liste d'instructions}, une commande LaTeX prenant exactement un
argument.

\bigskip
Voici un exemple qui change le style utilisé pour le nom d'une fonction Python, au moment de sa définition
(c'est-à-dire après le mot-clé |def|).

\begin{Verbatim}
\SetPitonStyle
  { Name.Function = \bfseries \setlength{\fboxsep}{1pt}\colorbox{yellow!50} }
\end{Verbatim}

Dans cet exemple, |\colorbox{yellow!50}| doit être considéré comme le nom d'une fonction LaTeX qui prend exactement
un argument, puisque, habituellement, elle est utilisée avec la syntaxe |\colorbox{yellow!50}{text}|.

\medskip
\begingroup
\SetPitonStyle 
  { Name.Function = \bfseries \setlength{\fboxsep}{1pt}\colorbox{yellow!50} }
Avec ce réglage, on obtient : \piton{def cube(x) : return x * x * x }
\endgroup



\medskip
Les différents styles sont décrits dans la table \ref{Semantic}.

\begin{table}[htb]
\centering
\caption{Usage des différents styles}
\label{Semantic}
\begin{tabularx}{0.9\textwidth}{@{}>{\ttfamily}l>{\raggedright\arraybackslash}X@{}}
\toprule
\normalfont Style  & Usage \\
\midrule
Number & les nombres \\
String.Short & les chaînes de caractères courtes (entre \texttt{'} ou \texttt{"}) \\
String.Long & les chaînes de caractères longues (entre \texttt{'''} ou \texttt{"""}) sauf les chaînes de
              documentation \\
String & cette clé fixe à la fois |String.Short| et |String.Long| \\
String.Doc & les chaînes de documentation \\
String.Interpol & les éléments syntaxiques des champs des f-strings (c'est-à-dire les caractères \texttt{\{}, \texttt{\}} et \texttt{:}) \\
Operator & les opérateurs suivants : \texttt{!= == << >> - \~{} + / * \% = < > \& .} \verb+|+ \verb|@| \\
Operator.Word & les opérateurs suivants : |in|, |is|, |and|, |or| et |not| \\
Name.Builtin & la plupart des fonctions prédéfinies par Python \\
Name.Function & le nom des fonctions définies par l'utilisateur \emph{au moment de leur définition}, c'est-à-dire 
                après le mot-clé \verb|def| \\ 
Name.Decorator & les décorateurs (instructions débutant par \verb|@| dans les classes) \\
Name.Namespace & le nom des modules (= bibliothèques extérieures) \\
Name.Class & le nom des classes au moment de leur définition \\
Exception & le nom des exceptions prédéfinies (eg: SyntaxError) \\
Comment & les commentaires commençant par \texttt{\#} \\
Comment.LaTeX & les commentaires commençant par \texttt{\#\#} qui sont composés en LaTeX par \pkg{piton}
                (\texttt{\#\#} est une séquence d'échappement vers LaTeX) \\
Keyword.Constant & |True|, |False| et |None| \\
Keyword & les mots-clés suivants :
          \ttfamily assert, break, case, continue, del,
          elif, else, except, exec, finally, for, from, 
          global, if, import, lambda, non local,
          pass, raise, return, try, while,
          with, yield, yield from.\\
  \bottomrule
\end{tabularx}
\end{table}



\subsection{Définition de nouveaux environnements}

Comme l'environnement |{Piton}| a besoin d'absorber son contenu d'une manière spéciale (à peu près comme du texte
verbatim), il n'est pas possible de définir de nouveaux environnements directement au-dessus de l'environnement
|{Piton}|.

C'est pourquoi \pkg{piton} propose une commande |\NewPitonEnvironment|. Cette commande a la même syntaxe que la
commande classique |\NewDocumentEnvironment|.

\bigskip
Par exemple, avec l'instruction suivante, un nouvel environnement |{Python}| sera défini avec le même comportement
que l'environnement |{Piton}|:

\verb|\NewPitonEnvironment{Python}{}{}{}|

\bigskip
Si on souhaite un environnement |{Python}| qui compose le code inclus dans une boîte de \pkg{tcolorbox}, on peut
écrire:

\begin{verbatim}
\NewPitonEnvironment{Python}{}
  {\begin{tcolorbox}}
  {\end{tcolorbox}}
\end{verbatim}


\end{document}